\chapter{Feasibility Spike}

\section{Purpose}

The feasibility spike focuses on the riskiest architectural assumption in this checkpoint: that
nopCommerce can remain the commerce core while durable omnichannel integration is added around
it. The spike therefore asks whether the current codebase has credible seams for the first
target evolution step: recording an order, creating durable integration work, and processing
that work outside the customer checkout path.

This chapter is intentionally scoped. It does not claim that the outbox, Integration Worker,
RabbitMQ queues, warehouse adapter, retry policy, dead-letter queue, or support dashboard are
already implemented. Those elements are part of the proposed target architecture. The spike
checks whether they can be introduced without a rewrite, and whether the nopCommerce source
already contains the specific seams that ADR~2 and the evolution roadmap depend on.

\section{What Was Tested}

The spike was conducted by reading the nopCommerce source directly against the classes named
in the ADRs, rather than building a full runtime implementation. This was appropriate for the
checkpoint because the main risk was architectural fit, not runtime behaviour: the question was
whether nopCommerce already exposes a stable order-placement seam and a background execution
mechanism credible enough to support the proposed asynchronous integration path.

\begingroup
\small
\setlength{\tabcolsep}{4pt}
\renewcommand{\arraystretch}{1.15}
\begin{center}
\begin{tabular}{|p{0.28\textwidth}|p{0.35\textwidth}|p{0.27\textwidth}|}
\hline
\textbf{Spike question} & \textbf{Inspection performed} & \textbf{Result} \\
\hline
Can order creation be used as the trigger for external fulfillment work? &
Inspected \texttt{OrderProcessingService} around the successful order-placement path. &
Yes. The order is persisted, order items are created, notifications are sent, and an
\texttt{OrderPlacedEvent} is published after order creation. \\
\hline
Is the current event mechanism sufficient for external integration? &
Inspected \texttt{EventPublisher}. &
No. Consumers run in process; exceptions are logged; there is no durable retry, replay, or
dead-letter lifecycle. \\
\hline
Is there a background execution seam for retry or polling work? &
Inspected \texttt{ScheduleTaskRunner} and scheduled-task behaviour. &
Yes, with limits. Scheduled tasks can execute background work with locking and interval checks,
but they need explicit integration state to be reliable. \\
\hline
Can existing inventory and shipment concepts support Scenario C projections? &
Inspected warehouse inventory, stock history, product inventory adjustment, and shipment classes. &
Partially. The baseline has stock, reserved quantity, stock history, tracking, shipped,
delivered, and ready-for-pickup fields, but not source freshness, conflict flags, or
integration-state metadata. \\
\hline
\end{tabular}
\end{center}
\endgroup

\section{What Worked}

The order-placement seam sits at exactly the right level. It fires after the database commit,
which is where ADR~2's atomicity guarantee needs to be anchored: an outbox record can be
written in the same transaction without touching the checkout flow or the payment handling path.

The locking and skip-if-running behaviour of \texttt{ScheduleTaskRunner} is exactly what an
outbox publisher or retry task requires to avoid overlapping execution. A production Integration
Worker would run as a separately deployable process, but a scheduled task is architecturally
sufficient to prove the tactic without introducing a new deployment unit prematurely.

The gap in the existing stock and shipment records is not structural but additive. The commerce
model already carries the right quantity and lifecycle fields; what is missing is integration
metadata --- a source identifier, freshness timestamp, stale flag, conflict flag, and retry
count. Adding those fields extends the existing model rather than replacing it, which confirms
the data ownership approach described in Chapter~6.

\section{What Did Not Work Yet}

The current implementation does not contain a transactional outbox, RabbitMQ integration,
warehouse adapter, shipping adapter, retry scheduler, dead-letter table, or operator-facing
integration dashboard. Therefore, the current codebase cannot yet demonstrate the mandatory
Scenario C pressure point at runtime. If the warehouse is unavailable today, there is no
durable fulfillment command that remains pending and retries automatically after recovery.

The in-process event mechanism identified in the spike as insufficient for durable external
integration is not a deficiency in nopCommerce's design for its intended purpose. It is a
structural gap relative to Scenario C, and it is precisely what ADR~2 addresses. The spike
therefore confirms the need for the outbox rather than raising doubt about the approach.

\section{Evidence}

The three most critical findings are supported by direct source references from the nopCommerce
repository.

\textbf{Insertion point in \texttt{OrderProcessingService}.}
The \texttt{PlaceOrderAsync} method in
\texttt{Nop.Services/Orders/OrderProcessingService.cs} ends its successful path with the
following sequence (lines 1608--1621):

\begin{verbatim}
    //notifications
    await SendNotificationsAndSaveNotesAsync(order);
    //reset checkout data
    await _customerService.ResetCheckoutDataAsync(...);
    await _customerActivityService.InsertActivityAsync(...);
    //raise event
    await _eventPublisher.PublishAsync(new OrderPlacedEvent(order));
    //check order status
    await CheckOrderStatusAsync(order);
\end{verbatim}

The order and its line items are already committed to the database before this point. An outbox
record inserted here --- in the same transaction scope --- would satisfy the atomicity guarantee
required by ADR~2 without touching the checkout flow or the payment handling path.

\textbf{Durability gap in \texttt{EventPublisher}.}
The \texttt{PublishAsync} method in \texttt{Nop.Services/Events/EventPublisher.cs} resolves
registered consumers and dispatches each one in sequence:

\begin{verbatim}
    foreach (var consumer in consumers)
    {
        try { await consumer.HandleEventAsync(@event); }
        catch (Exception exception)
        {
            await logger.ErrorAsync(exception.Message, exception);
        }
    }
\end{verbatim}

Exceptions are caught, logged, and swallowed. No failed work is persisted as a retryable
record. This is the structural gap that ADR~2 exists to close.

\textbf{Locking behaviour in \texttt{ScheduleTaskRunner}.}
The \texttt{ExecuteAsync} method in \texttt{Nop.Services/ScheduleTasks/ScheduleTaskRunner.cs}
skips execution if a previous run is still in progress and otherwise acquires a distributed
lock before executing:

\begin{verbatim}
    if (IsTaskAlreadyRunning(scheduleTask))
        return;
    ...
    await _locker.PerformActionWithLockAsync(
        scheduleTask.Type, expiration,
        () => PerformTaskAsync(scheduleTask));
\end{verbatim}

This guarantees that an outbox publisher or retry task does not overlap with a previous
execution, which is the minimum safety property required for the outbox publishing loop.

The remaining findings --- the inventory and shipment fields available on
\texttt{ProductWarehouseInventory}, \texttt{StockQuantityHistory}, and \texttt{Shipment} ---
were confirmed by inspecting the corresponding entity classes. They confirm that the extension
model described in Chapter~6 is additive rather than structural: the commerce records are
already there, and only integration metadata fields are missing.
\section{Draft Implementation Scope}

The implementation spike focuses on the smallest end-to-end slice that proves reliable asynchronous fulfillment propagation from nopCommerce to an external warehouse dependency.

\begin{enumerate}
    \item Add an \texttt{OmnichannelIntegrationTask} or outbox table with order ID, message
    type, payload, status, retry count, next attempt time, correlation ID, idempotency key,
    and last error.
    \item Create one outbox record when an order is accepted. The record should represent
    \texttt{FulfillmentRequested}.
    \item Add a scheduled task that reads pending records and calls a warehouse simulator or
    stub endpoint.
    \item Make the simulator return success, timeout, and failure modes through configuration.
    \item On failure, update the task to retrying or failed with a visible error. On success,
    mark it accepted or recovered.
    \item Add a minimal admin or log-based view showing pending, retrying, failed, and
    recovered states for one order.
\end{enumerate}

This slice directly demonstrates QA Scenario 1: checkout records the order, warehouse failure
does not block the customer path, the fulfillment request is retained as an observable record,
and recovery is visible after the dependency becomes available again. Shipping, inventory
conflict handling, and dead-letter intervention remain documented target behaviour, while this
spike confirms that the same architectural seams can support them.
